\documentclass[11pt]{article}
\usepackage[T1]{fontenc}
\usepackage{textcomp}
\usepackage[margin=0.75in]{geometry}
\usepackage{amsmath,amssymb,amsthm}
\usepackage{array,booktabs}
\usepackage{hyperref}
\setlength{\parindent}{0pt}
\newtheorem{theorem}{Theorem}
\theoremstyle{definition}
\newtheorem{definition}{Definition}
\title{Flow Proof Artifact: Nat}
\author{Generated by \texttt{flow doc proof}}
\date{Flow Proof Book}
\begin{document}
\maketitle
Peano definitions for addition on natural numbers.\\[0.5em]
\textbf{Source.} peano — https://en.wikipedia.org/wiki/Peano\_axioms\\[0.5em]
\bigskip
\noindent{\large \textbf{Definition 1.} \textit{Adding zero on the left gives the other number} \hfill the natural numbers \cdot addition \cdot zero is the left identity}\par
\smallskip\noindent\textit{Coordinate: the natural numbers \cdot addition \cdot zero is the left identity}\par
\smallskip\noindent\textit{Source: peano}\par
\medskip
\noindent\textbf{Goal.} Adding zero on the left gives the other number.  (0 + 7 = 7)\par
\noindent$\displaystyle \forall m \in \mathbb{N}\quad 0 + m = m$\par
\medskip
\noindent
\begin{tabular}{@{} >{\raggedright\arraybackslash}p{0.06\textwidth} >{\raggedright\arraybackslash}p{0.47\textwidth} >{\raggedleft\arraybackslash}p{0.41\textwidth} @{}}
\textbf{\#} & \textbf{Proof} & \textbf{Mathematics} \\
\midrule
\textbf{1.} & We stipulate zero is the left identity for addition on the natural numbers — this is a definition, not a derived fact. & \\[0.45em]
\textbf{2.} & This follows directly from the definition: 0 plus m equals m. Hence proven. & $\displaystyle 0 + m = m$ \\[0.45em]
\bottomrule
\end{tabular}
\label{thm:Nat-addition-zero_is_the_left_identity}
\par\medskip\hrule
\bigskip
\noindent{\large \textbf{Definition 2.} \textit{Adding one more on the right steps the sum by one} \hfill the natural numbers \cdot addition \cdot successor on the right steps the sum}\par
\smallskip\noindent\textit{Coordinate: the natural numbers \cdot addition \cdot successor on the right steps the sum}\par
\smallskip\noindent\textit{Source: peano}\par
\medskip
\noindent\textbf{Goal.} Adding one more on the right steps the sum by one\par
\noindent$\displaystyle \forall n \in \mathbb{N} \forall m \in \mathbb{N}\quad n + \mathrm{succ}(m) = \mathrm{succ}(n + m)$\par
\medskip
\noindent
\begin{tabular}{@{} >{\raggedright\arraybackslash}p{0.06\textwidth} >{\raggedright\arraybackslash}p{0.47\textwidth} >{\raggedleft\arraybackslash}p{0.41\textwidth} @{}}
\textbf{\#} & \textbf{Proof} & \textbf{Mathematics} \\
\midrule
\textbf{1.} & We stipulate successor on the right steps the sum for addition on the natural numbers — this is a definition, not a derived fact. & \\[0.45em]
\textbf{2.} & This follows directly from the definition: n plus the successor of m equals the successor of n plus m. Hence proven. & $\displaystyle n + \mathrm{succ}(m) = \mathrm{succ}(n + m)$ \\[0.45em]
\bottomrule
\end{tabular}
\label{thm:Nat-addition-successor_on_the_right_steps_the_sum}
\par\medskip\hrule
\end{document}
